\section{Positive slopes and the six-contact product}

Let
\[
 \Theta=\{0,-1/2,-1/4,-2/5,-1,-1/3\},\qquad
 F(t)=\prod_{\theta\in\Theta}(t-\theta)
 =\frac{t(t+1)(2t+1)(3t+1)(4t+1)(5t+2)}{120},
 \qquad R(t)=\frac{F(t)}{(t-1)^8}.
\]
For a nonzero rational number $x$ and a finite prime set $S$, let
$\operatorname{Odd}_S(x)$ be the product of primes outside $S$
whose valuation in $x$ is odd. Negative valuations are included.

\begin{proposition}
The rational involution
\[
                      \tau(t)=\frac{t+1}{2t-1}
\]
preserves $\Theta$ and satisfies
\[
 F(\tau(t))=\frac{27F(t)}{(2t-1)^6},\qquad
 R(t)R(\tau(t))
 =3\left(\frac{3F(t)(2t-1)}{(t-1)^4(t-2)^4}\right)^2.
\]
For every integer $m\ge2$, put
\[
 \lambda_m=\frac{m^2+1}{2m^2},\qquad
 \mu_m=\frac{3m^2+1}{2},\qquad
 a_m=\frac{m^2(3m^2+1)}{m^2+1}.
\]
Then $\mu_m=a_m\lambda_m=\tau(\lambda_m)$, $a_m>4$,
$\lambda_m,\mu_m>0$, and neither argument is a contact or pole.
Moreover $\lambda_m$ is a sum of two rational squares and
$\operatorname{Odd}_S(R(\lambda_m)R(a_m\lambda_m))=1$
whenever $3\in S$, whereas $h(a_m)\to\infty$.
\end{proposition}
\begin{proof}
The first identity is obtained by substitution into the six factors;
the second uses $\tau(t)-1=(2-t)/(2t-1)$. Direct calculation gives
all formulas for $\lambda_m,\mu_m,a_m$. The norm representation is
\[
 \lambda_m=\left(\frac{m+1}{2m}\right)^2+
            \left(\frac{m-1}{2m}\right)^2.
\]
Since $1/2<\lambda_m\le5/8$ and $\mu_m\ge13/2$, the
arguments avoid the contacts, $1$, and $2$. The product is three
times a nonzero rational square, so its odd radical outside $S$
is one. Finally $a_m>3m^2-3$, whence $h(a_m)\ge\log a_m\to\infty$.
\end{proof}

Thus requiring a positive slope, even with an explicit
$\mathbb Q(i)$-norm representation, does not by itself yield any
bound for the scale in terms of the odd radical of the product.
The proposition does not assert that $a_m$ has the further form
$\beta^8+T(n)^8$, or that it extends to a solution of the complete
Diophantine source.

\begin{proposition}
The following table lists every rational M\"obius transformation
preserving $\Theta$. A row $(a,b,c,d;k)$ means
\[
 \gamma(t)=\frac{at+b}{ct+d},\qquad
              F(\gamma(t))(ct+d)^6=kF(t).
\]
\[
\begin{array}{rrrr|r}
a&b&c&d&k\\\hline
1&0&-5&-1&-1\\
1&0&0&1&1\\
1&1&-7&-4&-27\\
1&1&2&-1&27\\
2&1&-7&-3&1\\
2&1&-3&-2&-1\\
3&1&-8&-3&-1\\
3&1&-7&-2&1\\
4&1&-13&-4&27\\
4&1&-7&-1&-27\\
5&2&-14&-5&-27\\
5&2&-11&-5&27
\end{array}
\]
Apart from the identity, $\tau$ is the only map in the table
that sends some positive rational number to a positive number.
Its positive domain is $t>1/2$. If $t>1$, then
$0<\tau(t)/t<2$. Consequently no nonidentity contact automorphism
can satisfy $\gamma(\lambda)=a\lambda$ with $\lambda>1$ and $a>4$.
\end{proposition}
\begin{proof}
A M\"obius transformation is determined by the images of the
three distinct points $0,-1/2,-1/4$. Sending these to each of the
$6\cdot5\cdot4=120$ ordered triples in $\Theta$, and then checking
the remaining three contacts, gives exactly the twelve displayed
matrices. The calculation uses only rational $3\times3$
determinants; the accompanying exact checker supplies the resulting
permutations and polynomial identities. Alternatively each of the
120 candidates can be checked directly from the displayed defining
rule. The positivity assertions follow from the signs of the matrix
entries. For $t>1$, the inequality $\tau(t)/t<2$ is equivalent
to $4t^2-3t-1=(t-1)(4t+1)>0$.
\end{proof}

Requiring $\lambda>1$ therefore removes these automorphism
families. It does not prove a height bound for the product: other
pairs on quadratic twists of $y^2=F(t)$ must still be controlled.
Likewise, the identity
$\tau(\lambda)=a\lambda$ has a rational solution precisely when
$a^2+10a+1$ is a rational square. No conclusion about its
intersection with the prescribed eighth-power scales is claimed.
